PP computation is implemented in \texttt{crates/bisect-analysis/src/compactness.rs}.
The implementation uses the \texttt{geo} crate for polygon area and perimeter
computation and \texttt{proj} for coordinate system transformation.

\subsection{Computational Pipeline}

\begin{enumerate}
  \item \textbf{Load geometries}: District polygons are loaded from the run's
    output shapefile in WGS84 (EPSG:4326).
  \item \textbf{Project}: Each district polygon is projected to Albers Equal Area
    Conic (EPSG:5070) using the PROJ library.
  \item \textbf{Compute area}: Shoelace formula applied to the projected polygon
    vertices gives $A$ in square metres.
  \item \textbf{Compute perimeter}: Euclidean distance summed over projected edge
    segments gives $P$ in metres.
  \item \textbf{Apply formula}: PP $= 4\pi A / P^2$.
\end{enumerate}

\subsection{Numerical Stability}

For very small districts (single-tract districts in single-state analyses),
$P^2$ can be large relative to $A$, causing PP to approach zero.
No numerical instability occurs because PP is bounded below by zero and
the formula involves only division, not subtraction of nearly equal quantities.

\subsection{CLI Usage}

PP is computed by:
\begin{verbatim}
bisect label-analyze <label> --year 2020 --types pp
\end{verbatim}
or as part of the full compactness profile:
\begin{verbatim}
bisect label-analyze <label> --year 2020 --types all
\end{verbatim}
Output JSON format:
\begin{verbatim}
{
  "districts": [
    {
      "district_id": 1,
      "polsby_popper": 0.231,
      "projection": "EPSG:5070"
    },
    ...
  ],
  "state_mean_pp": 0.224,
  "seed": 42
}
\end{verbatim}

\subsection{Test Suite}

L0 invariants (inline unit tests):
\begin{itemize}
  \item PP of a circle approximated by a regular 1000-gon is within $0.001$ of $1.0$.
  \item PP is in $[0,1]$ for any valid polygon.
  \item PP of a square is within $0.001$ of $\pi/4 \approx 0.785$.
  \item PP is scale-invariant: doubling all coordinates changes PP by $< 10^{-9}$.
\end{itemize}

L1 integration test:
\begin{itemize}
  \item On a $4 \times 4$ grid of unit squares with two equal rectangular districts,
    PP is computed without panic and matches the formula value within $0.001$.
\end{itemize}

L2 regression test:
\begin{itemize}
  \item NC 2020 ratio-optimal mean PP $\geq 0.22$ (seed $= 42^\dagger$).
  \item Projection to EPSG:5070 increases mean PP by $\geq 0.01$ for NC coastal districts.
\end{itemize}
